\subsection{Fredholm 指标}\label{sec:fredholm-index}

对于许多线性 PDE, 证明唯一性远比证明存在性容易. 众所周知, 对于有限维向量空间中的算子, 唯一性与存在性事实上等价; 这称为 Fredholm 二择一结论. 研究无穷维情形下满足 Fredholm 二择一结论的算子很重要. 先给出一个定义.

\begin{Definition}
\label{def:semi-fredholm-and-fredholm-operators}
设 $X$ 和 $Y$ 为 Banach 空间. 若算子 $A\in\mathcal{L}(X,Y)$ 的值域 $\mathcal{R}(A)$ 闭, 且 $\mathcal{N}(A)$ 有限维或 $\mathcal{R}(A)$ 在 $Y$ 中有有限维补空间, 则称该算子为\textbf{半 Fredholm 算子}. 若两者都成立, 则称该算子为 \textbf{Fredholm 算子}.
\end{Definition}

上面的定义限于有界算子. 然而, 若 $A$ 无界, 则总可以将其看作定义在 Banach 空间 $\mathcal{D}(A)$ 上的有界算子, 其中 $\mathcal{D}(A)$ 赋予图范数 (参见 Problem~\ref{ex:operator-graph-norm-completeness}).

半 Fredholm 算子的最重要性质是一个称为 Fredholm 指标的量.

\begin{Definition}
\label{def:fredholm-nullity-deficiency-index}
设 $A\in\mathcal{L}(X,Y)$ 为半 Fredholm 算子. 称 $\mathcal{N}(A)$ 的维数为 $A$ 的\textbf{零度}, 称 $\mathcal{R}(A)$ 的补空间的维数为 $A$ 的\textbf{亏数} (若 $\mathcal{R}(A)$ 没有有限维补空间, 则亏数为无穷). 量
\begin{equation}\label{eq:fredholm-index-nullity-minus-deficiency}
\operatorname{ind}A=\operatorname{nul}A-\operatorname{def}A
\end{equation}
称为 $A$ 的 \textbf{(Fredholm) 指标}.
\end{Definition}

若 $A$ 为 Fredholm 算子, 则指标有限; 否则为正无穷或负无穷. 关于半 Fredholm 算子的关键定理如下:

\begin{Theorem}
\label{thm:semi-fredholm-index-local-stability}
设 $A\in\mathcal{L}(X,Y)$ 为半 Fredholm 算子. 那么存在 $\varepsilon>0$, 使得任意满足 $\|B-A\|<\varepsilon$ 的 $B\in\mathcal{L}(X,Y)$ 也是半 Fredholm 算子, 且 $\operatorname{ind}B=\operatorname{ind}A$.
\end{Theorem}

下面给出的证明适用于 $A$ 为 Fredholm 算子, 或 $X$ 和 $Y$ 为 Hilbert 空间的情形. 一般情形的困难在于, Banach 空间的闭子空间未必有闭补空间; 一般情形的证明参见 \cite{Ka}. 对于 $A$ 为 Fredholm 算子的情形, 注意以下引理:

\begin{Lemma}
\label{lem:finite-dimension-or-codimension-closed-complement}
设 $X$ 为 Banach 空间, $V$ 为 $X$ 的闭子空间, 且有限维或具有有限余维. 那么存在 $X$ 的闭子空间 $W$, 使得 $X=V\oplus W$.
\end{Lemma}

\begin{Proof}
若 $V$ 具有有限余维, 只需注意到每个有限维赋范向量空间都完备, 因而 $X$ 的每个有限维子空间都闭. 若 $V$ 有限维, 设 $e_i$, $i=1,\ldots,n$, 为 $V$ 的一组基. 由 Hahn-Banach 定理, 可以构造线性泛函 $x_i^*\in X^*$, 使得 $x_i^*(e_j)=\delta_{ij}$. 然后将 $W$ 定义为各个 $x_i^*$ 的零空间之交.
\end{Proof}

在下面的证明中, 还要用到闭子空间与有限维子空间的直和是闭的这一事实; 其证明留作练习 (Problem~\ref{ex:closed-plus-finite-dimensional-subspace}).

现在假设 $A$ 为 Fredholm 算子, 或 $X$ 和 $Y$ 为 Hilbert 空间, 继续证明该定理.

\begin{Proof}
设 $V$ 和 $W$ 分别为 $X$ 和 $Y$ 的子空间, 使得 $X=\mathcal{N}(A)\oplus V$, $Y=\mathcal{R}(A)\oplus W$. 用 $\widetilde A(v,w)=Av+w$ 定义 $\widetilde A:V\times W\to Y$. 那么 $\widetilde A$ 是双射. 对某个给定算子 $B\in\mathcal{L}(X,Y)$, 类似地定义 $\widetilde B$. 若 $\|B-A\|$ 足够小, 则 $\|\widetilde B-\widetilde A\|$ 也足够小, 因而 $\widetilde B$ 是双射. 换言之, 对给定的 $y\in Y$, 方程 $Bv+w=y$ 有唯一解 $v\in V$, $w\in W$. 立即得到 $\mathcal{R}(B)+W=Y$, 即 $\operatorname{def}B\leq\operatorname{def}A$, 以及 $\mathcal{N}(B)\cap V=\{0\}$, 即 $\operatorname{nul}B\leq\operatorname{nul}A$. 此外, $B(V)=\{y\in Y\mid\widetilde B^{-1}(y)\in V\}$ 是闭的, 且在 $\mathcal{R}(B)$ 中具有有限余维, 因为 $B(V)$ 在 $Y$ 中具有有限余维, 或 $V$ 在 $X$ 中具有有限余维. 因而 $\mathcal{R}(B)$ 是闭的, $B$ 是半 Fredholm 算子.

由于 $V$ 在 $X$ 中具有有限余维, 或 $B(V)$ 在 $Y$ 中具有有限余维, 得到 $V+\mathcal{N}(B)$ 在 $X$ 中具有有限余维. 现在设 $Z$ 为满足 $X=V\oplus\mathcal{N}(B)\oplus Z$ 的空间. 那么有 $Y=B(V)\oplus W$ 和 $\mathcal{R}(B)=B(V)\oplus B(Z)$, 即 $\operatorname{def}B=\dim W-\dim B(Z)=\dim W-\dim Z$ 以及 $\operatorname{nul}B=\dim\mathcal{N}(A)-\dim Z$; 两个指标相等立即得证.
\end{Proof}

由 Theorem~\ref{thm:semi-fredholm-index-local-stability} 得到以下推论:

\begin{Corollary}
\label{cor:semi-fredholm-open-set-constant-index}
所有半 Fredholm 算子组成的集合在 $\mathcal{L}(X,Y)$ 中是开集. 此外, 指标在每个连通分支上为常数.
\end{Corollary}

因此, 若 $A(t)\in\mathcal{L}(X,Y)$ 对每个 $t$ 均为半 Fredholm 算子, 且在 $t\in[a,b]$ 上连续依赖于 $t$, 则指标与 $t$ 无关. 这在应用中常常很有用, 因为 $A(a)$ 的指标可能比 $A(b)$ 的指标更容易求得. 以下定理就是这种方法的一个应用.
